% =====================================================================
%  27.5  Recognition and the method boundary -- the Mo signature
%  (offline + overlapping + non-invertible), and the six hand-offs:
%  invertible -> Euler+Fenwick (Ch.5); online path -> HLD (Ch.22) / LCT
%  (Ch.25); per-set fold -> virtual tree (Ch.26); small marked sets /
%  subtree folds -> DSU-on-tree (Ch.17); all-pairs distance shapes ->
%  centroid (Ch.24); LCA-only batches -> Tarjan (this chapter, 27.4).
%  WORKED pick examples only -- no solvedproblem (contract: two, spent).
%  Exports: sec:c21-boundary  def:c21-signature  def:c21-pick  lst:c21-kit
%  CF 86D (Powerful array) and CF 375D are PROSE references only -- their
%  URLs live in other chapters; nothing here links a foreign URL.
% =====================================================================
\section{The Method Boundary: When Mo Is Wrong}\label{sec:c21-boundary}

\begin{previouslybox}[Ch.~26 ended at a boundary -- it was this one]
The last chapter of this book's previous stop, virtual trees (\chref{c30}),
closed its pick table with an honest confession: "many overlapping offline
node-set queries whose statistic will not fold per-set" does not compress onto
any virtual tree, and belongs to Mo.  You now own Mo -- so the debt must be
collected in the other direction: \emph{which shapes do NOT belong to it}?  A
tool this comfortable (\(\bigO{(n+q)\sqrt n}\) answers from one sweep) is also
this easy to misfire: onto online queries (block order dies), onto invertible
statistics (a BIT is \(\sqrt n/\log n\) times cheaper), onto per-set folds a
virtual tree eats in \(\bigO{k\log k}\).  This section is the boundary contract
for the whole chapter.
\end{previouslybox}

\begin{topicintro}
Three questions, asked in order, decide everything.  \emph{Are all queries
known before any answer is needed?}  No \(\Rightarrow\) Mo is illegal -- online
path shapes go to HLD (\chref{c16}) or the link-cut tree (\chref{c23}).
\emph{Does the statistic telescope over a prefix?}  Yes \(\Rightarrow\) Mo is
overkill -- Euler flatten + Fenwick/segment answers in
\(\bigO{(n+q)\log n}\) (Ch.~18's CF 383C chose this door, \(\chref{c14}\); the
root-path batch CF 463E in \(\secref{c21-subtree}\) chose the stack instead).
\emph{Are the queried sets subtrees/paths of ONE fixed tree, many of them
overlapping heavily?}  Yes \(\Rightarrow\) this chapter's engine.  The three
negative escapes -- small marked sets (DSU-on-tree, \chref{c28}), per-set
compression (virtual tree, \chref{c30}), distance distributions over all pairs
(centroid, \chref{c18}) -- are worked through below, then \defref{c21-pick}
codifies the order of the questions.

\textbf{Prerequisites.} The whole chapter, plus the boundary sections of
Ch.~17/22/24/26 it hands work back to.

\textbf{After this section you can}
\begin{itemize}[nosep]
  \item recite the Mo signature (offline, overlapping, non-invertible, \(O(1)\) update); \hfill\lv{4}
  \item reject Mo correctly for online, invertible, and per-set-fold shapes; \hfill\lv{5}
  \item choose Mo vs.\ DSU-on-tree for subtree statistics with a reason; \hfill\lv{5}
  \item price each alternative before committing \(\bigO{(n+q)\sqrt n}\). \hfill\lv{4}
\end{itemize}
\end{topicintro}

% ---------------------------------------------------------------
\subsection{The signature}
% ---------------------------------------------------------------

\begin{definition}{The Mo signature (four tests)}{c21-signature}
A problem is Mo-shaped iff all four hold:
\begin{enumerate}[nosep,leftmargin=1.6em]
  \item \textbf{Offline:} every query is known upfront; no query's answer
        feeds an update (reordering the questions must be legal).
  \item \textbf{Overlapping ranges:} the query sets are ranges/paths of
        \emph{one fixed tree} that share elements heavily -- otherwise there
        is no work to amortise.
  \item \textbf{Non-invertible statistic:} distinct counts,
        \(\sum\mathrm{freq}^2\), mode, "how many values \(\ge l\) times" --
        no prefix/subtraction form exists.
  \item \textbf{\(O(1)\) (de)activation:} adding or removing \emph{one node}
        updates the statistic in \(O(1)\) (or \(O(\log)\) with a small
        structure -- then re-price).
\end{enumerate}
Fail any test and the table below names where the problem really lives.
\end{definition}

\begin{recognitionbox}[Question shapes: Mo versus each look-alike]
\begin{itemize}[nosep,leftmargin=1.4em]
  \item \textbf{"distinct colours on subtree\((v)\)" for many \(v\), all queries
        given at once.}
        \reckey{tests 1--4 all pass: flatten + path-free Mo.}
        \recuse{\item the 27.2 transfer: window over \([\tin(v),\tout(v)]\).}
        \recnotuse{\item DSU-on-tree also solves it (\chref{c28}) -- see the pick
        example below; price both, then move.}
  \item \textbf{"sum of values on subtree\((v)\)" with point updates
        interleaved.}
        \reckey{invertible \emph{and} online: Euler + Fenwick/segment,
        \(\bigO{\log n}\) per op (\chref{c24}).}
        \recuse{\item the difference/prefix form.}
        \recnotuse{\item Mo cannot see query \(i{+}1\) after an update at \(i\).}
  \item \textbf{"path \(u\)-\(v\) aggregate, answers needed online, edge updates
        allowed".}
        \reckey{the HLD contract (\chref{c16}); with topology-like operations,
        the LCT (\chref{c23}).}
        \recuse{\item decompose path \(\to\) \(O(\log n)\) segment ranges.}
        \recnotuse{\item Mo: test 1 fails immediately.}
  \item \textbf{Same path statistic but offline.}
        \reckey{path Mo (\(\secref{c21-path}\)): doubled tour + toggle.}
        \recuse{\item block order on tour windows.}
        \recnotuse{\item HLD when \(q\le 2\cdot10^5\) and updates exist -- Mo has
        lost test 1 again.}
  \item \textbf{"aggregate over each query's small marked set" (\(k\) marked,
        \(\sum k\) moderate).}
        \reckey{virtual tree (\chref{c30}): build \(\bigO{k\log k}\), fold.}
        \recuse{\item the sets are \emph{different} per query and do not overlap
        as ranges.}
        \recnotuse{\item Mo -- windows would shuffle arbitrary non-controlling
        node sets, no amortisation.}
  \item \textbf{"count pairs at distance \(d\)", "sum over all paths" (no query
        list at all).}
        \reckey{centroid decomposition (\chref{c18}) -- the object is the whole
        tree's path multiset.}
        \recuse{\item divide-and-conquer on paths.}
        \recnotuse{\item Mo -- there are no ranges to order.}
\end{itemize}
\end{recognitionbox}

% ---------------------------------------------------------------
\subsection{The pick, worked: Mo versus DSU-on-tree}
% ---------------------------------------------------------------

The sharpest boundary inside this chapter's own regime (offline, subtree,
non-invertible) is Mo vs.\ DSU-on-tree (\chref{c28}, small-to-large): both
answer "distinct colours in subtree\((v)\) for all \(v\)" -- CF 375D is the
textbook instance (its full treatment and URL live there; named in prose only
here, by contract).

\begin{center}\renewcommand{\arraystretch}{1.25}\setlength{\tabcolsep}{3pt}
\begin{tabular}{@{}p{0.20\linewidth} p{0.24\linewidth} p{0.24\linewidth} >{\raggedright\arraybackslash}p{0.18\linewidth}@{}}
\toprule
\textbf{Axis} & \textbf{Mo (this chapter)} & \textbf{DSU-on-tree (Ch.17)} & \textbf{Wins}\\
\midrule
Shape & any ranges/paths, \(q\) arbitrary & subtree queries, all \(v\) (or many) & Mo (generality) \\
Time & \(\bigO{(n+q)\sqrt n}\) & \(\bigO{n\log n}\) (subtrees only) & DSU (on subtrees) \\
Paths & native (27.3 toggle) & awkward (keeps subtrees, not paths) & Mo \\
Answers & stored offline, printed at end & available for every node at end & tie \\
Stat update & must be \(O(1)\) per move & must merge maps small-to-large & Mo (simpler contract) \\
Memory & \(\bigO{n+q}\) flat & \(\bigO{n}\) + per-node maps & Mo \\
\bottomrule
\end{tabular}
\end{center}

\begin{tracebox}[Worked pick: three statements, three verdicts]
\textbf{(a)} \(n,q\le10^5\); query: distinct \(\mathit{val}\) on path \(u\)-\(v\);
all queries given upfront.  Tests: 1 pass, 2 pass, 3 pass (distinct), 4 pass
(toggle) \(\Rightarrow\) \emph{path Mo}, \(\bigO{(n+q)\sqrt n}\) \(\approx\)
\(6\cdot10^7\) moves -- fits.  DSU-on-tree cannot hold a path.
\textbf{(b)} \(n\le10^5\); same statistic for \emph{every subtree} (query set
\(=\) all \(v\)).  Test 3 passes, but the subtree-all-queries specialisation is
\(\bigO{n\log n}\) with Ch.~17's keep-heavy-child trick: DSU-on-tree by
\(\sqrt n/\log\) factor.  Verdict \emph{DSU-on-tree} -- Mo still \emph{correct},
just slower; know the price you are paying.
\textbf{(c)} \(n,q\le2\cdot10^5\); subtree distinct, but queries arrive one per
line and an update \(a_x := y\) may follow each answer.  Test 1 fails
\(\Rightarrow\) neither Mo nor plain DSU: the honest door is sqrt-decomposition
\emph{over time} (rebuild DSU answers per block of operations) -- beyond this
chapter; recorded so the rejection names a destination, not a cliff. \(\Box\)
\end{tracebox}

\begin{definition}{The order of the four questions (pick protocol)}{c21-pick}
\emph{Online?} \(\Rightarrow\) HLD/LCT (\chref{c16}, \chref{c23}) or Euler+segment
for subtrees.  \emph{Invertible?} \(\Rightarrow\) flatten + Fenwick
(\chref{c24}); the parity-normalised CF 383C (Ch.~18, \(\chref{c14}\)) is the
template.  \emph{All query sets
small marked subgraphs?} \(\Rightarrow\) virtual tree (\chref{c30}).
\emph{Whole-tree path distribution?} \(\Rightarrow\) centroid (\chref{c18}).
\emph{Subtree-only, every node?} \(\Rightarrow\) price DSU-on-tree (\chref{c28})
first.  \emph{Still here?} The problem is Mo's by elimination -- and it is
\(\bigO{(n+q)\sqrt{n+q}}\) (array/subtree) or path Mo with the doubled tour.
\end{definition}

\begin{lstlisting}[caption={The pick protocol as a decision skeleton -- paste-ready comments for contest code; each branch names the chapter that owns the technique}, label={lst:c21-kit}]
// 1. do answers have to come BETWEEN updates?  -> NO Mo.
//      path ranges: HLD (Ch.22). dynamic topology: LCT (Ch.25).
// 2. can the statistic be written pref(r) - pref(l-1)? -> NO Mo.
//      subtree sums / parity-normalised adds: Euler + Fenwick (Ch.5).
// 3. is each query a SMALL marked set to fold over?     -> virtual tree (Ch.26)
// 4. no query list at all -- whole-tree pair counts?     -> centroid (Ch.24)
// 5. subtree-only, statistic mergeable small-to-large?   -> DSU on tree (Ch.17)
// 6. LCA batch only, offline, want linear-ish?           -> Tarjan (27.4)
// 7. otherwise: offline, overlapping ranges/paths, O(1) add/rm:
//        subtrees -> flatten P, Mo on [tin,tout]         (27.2)
//        paths    -> doubled tour E, toggle, lca split   (27.3)
//        B = (int)( (2*n) / sqrt(max(1,q)) ) ; odd-even R order.
\end{lstlisting}

\begin{insightbox}[Mo is the tool of last resort -- and that is a compliment]
Every rejection above is a statement that something \emph{cheaper} applies:
invertibility buys \(\log\), structure buys \(\log\), smallness buys \(k\log k\).
Mo is what remains when no structure exists -- but it always exists, needing
only test 4.  "I cannot compress it, so I will order it" is the mature
contest reflex: \(\bigO{(n+q)\sqrt n}\) with a 12-line engine beats a
\(\bigO{n\log^2 n}\) structure you cannot debug at minute 80.
\end{insightbox}

\begin{notebox}[Boundary homework this chapter leaves elsewhere]
Three shapes keep walking out of these doors and get settled by later
chapters: online path queries with updates arrive at HLD (\chref{c16}) and its
dynamic cousin (\chref{c23}); "tree isomorphism / canonical form" questions --
which subtree equals which -- belong to Chapter~28, next door, where hashing
supplies the equality test Mo never needed (Mo compares counts, not trees).
And if your "non-invertible statistic" is actually "frequency of frequencies"
\emph{on root paths}, the Summit at the end of this chapter answers it directly.
\end{notebox}

\begin{complexitybox}[The boundary, priced]
Same input class (offline subtree/path, non-invertible): Mo
\(\bigO{(n+q)\sqrt n}\) general \(\to\) DSU-on-tree \(\bigO{n\log n}\) (subtrees
only).  Invertible: \(\bigO{(n+q)\log n}\) Euler+Fenwick.  Online paths: HLD
\(\bigO{(n+q\log n)\log n}\)-ish with updates.  Per-set: virtual tree
\(\bigO{k\log k}\).  LCA batch: Tarjan \(\bigO{(n+q)\alpha}\).  The gaps are a
factor \(\sqrt n\) to \(\log^2 n\) -- enough to flip AC/TLE at \(n=2\cdot10^5\),
which is why the protocol order in \defref{c21-pick} is not negotiable.
\end{complexitybox}

\begin{recallbox}[27.5 in one breath]
\emph{Four tests make a Mo problem: offline, overlapping, non-invertible,
\(O(1)\) update (\(\defref{c21-signature}\)).  Each failed test is not a dead
end but a hand-off: online \(\to\) HLD/LCT; invertible \(\to\) Fenwick; small
marked \(\to\) virtual tree; all-pairs \(\to\) centroid; subtree-all \(\to\)
DSU-on-tree; LCA-only \(\to\) Tarjan.  Ask the questions in this order
(\(\defref{c21-pick}\), \(\coderef{c21-kit}\) -- and only then fire the
sweep.}
\selftest{
  \item Which single test does an interleaved-update problem fail, and what should you reach for on paths? \emph{(Test 1, offline; HLD (Ch.22), or LCT (Ch.25) for dynamic trees.)}
  \item Subtree distinct-colours for \emph{every} node: Mo or DSU-on-tree, and why? \emph{(DSU-on-tree -- the specialised O(n log n) beats O((n+q) sqrt n) when queries are exactly all subtrees.)}
  \item Name the CF 383C lesson in one sentence. \emph{(Scary updates can still be invertible after normalisation -- always look for the prefix before paying the sqrt tax; Ch.~18 solved it that way.)}
  \item CF 463E passed both the BIT and Mo doors yet was solved with neither. What did it walk through instead? \emph{(The third door: the recursion stack carries the root-path window; per-key push/undo table plus one rebuild per block of the \(\le 50\) updates.)}
}
\end{recallbox}
